% ============================================================
% v2 — Appendici A–B (snelle, fuori conteggio pagine)
% A: iperparametri e configurazione (tabellati anziché dump YAML)
% B: riproducibilità
% Valori verificati dagli snapshot in latex/configs/ (2026-07-14).
% ============================================================

\chapter{Osservazioni e reward per componente}
\label{app:obs-reward}

Le due tabelle che seguono danno il dettaglio per componente degli spazi
di osservazione (\cref{sec:osservazioni}) e della funzione di reward
(\cref{sec:reward}), la cui struttura e il cui razionale sono descritti
nel \cref{ch:piattaforma}. Tutti i valori sono verificati dal codice del
trunk di campagna.

\begin{table}[h]
  \centering\small
  \caption{Layout delle osservazioni. Le barre indicano costanti di
  normalizzazione; le feature dei vicini sono nel sistema di riferimento
  dell'ego. TTC $=$ time-to-collision.}
  \label{tab:obs}
  \begin{tabular}{@{}llp{8.0cm}@{}}
    \toprule
    Dim. & Gruppo & Contenuto \\
    \midrule
    \multicolumn{3}{@{}l}{\emph{Veicolo (47D)}} \\
    0--12  & Ego e percorso   & velocità (/30), accelerazione (/10), modulo (/120 km/h), heading; distanza (/\SI{50}{m}) e bearing al waypoint successivo; offset laterale in corsia; progresso continuo \\
    13--23 & Semaforo, vicini & flag e stato del semaforo; 3 veicoli più prossimi entro \SI{50}{m} (offset long.\ e lat.\ /50, velocità relativa /30) \\
    24--33 & Anteprima e geometria & sterzo precedente; waypoint $+1/+3/+5$; errore di heading; angolo curva imminente; errore laterale con segno \\
    34--43 & Pedoni e hazard  & 2 pedoni più prossimi, stesse 3 feature; rischio TTC e occupazione del corridoio, vs.\ veicoli e vs.\ pedoni \\
    44--46 & Progresso/tempo  & passi senza progresso (/300); flag anti-loop; tempo rimanente \\
    \midrule
    \multicolumn{3}{@{}l}{\emph{Pedone (26D)}} \\
    0--11  & Ego e waypoint   & posizione ($x,y$ /500, $z$ /50); velocità (/5); vettore frontale; distanza (/\SI{100}{m}) e bearing al waypoint; flag marciapiede \\
    12--17 & Veicoli vicini   & 2 più prossimi entro \SI{50}{m}: stesse 3 feature \\
    18--25 & Percorso e hazard & completamento; waypoint $+1/+3$ (/100); errore di heading; rischio TTC e occupazione vs.\ veicoli \\
    \bottomrule
  \end{tabular}
\end{table}

\begin{table}[h]
  \centering\small
  \caption{Componenti della reward. Le righe veicolo ``gated'' sono
  sospese durante l'attesa a un semaforo rosso/giallo; le righe marcate
  $^{\dagger}$ richiedono inoltre $\mathit{safe\_to\_push}$
  ($\mathit{hazard} < 0.85$) e allineamento al percorso $> 0.25$.}
  \label{tab:reward}
  \begin{tabular}{@{}p{3.2cm}p{5.2cm}p{4.4cm}@{}}
    \toprule
    Componente & Trigger & Valore \\
    \midrule
    \multicolumn{3}{@{}l}{\emph{Veicolo}} \\
    Waypoint raggiunto / saltato & per waypoint avanzato / saltato & $+100$ / $-10$ \\
    Progresso denso       & variazione della distanza dal waypoint & $+4.0$ per metro guadagnato \\
    Collisione            & primo contatto (terminale) & $-500$ \\
    Fuori corsia / eccesso di velocità & $>$ \SI{2}{m} dal centro corsia \ / \ $v > 50$ km/h & $-0.5$ per metro \ / \ $-0.10$ per km/h \\
    Inattività (gated)    & velocità $< 2.5$ km/h & $-0.15$ \\
    Bonus di partenza (gated$^{\dagger}$) & acceleratore da fermo & $(0.20 + 0.30\,u)\cdot\text{thr}\cdot\text{align}$, freno simmetrico; $u = \min(\text{no-progress}/200, 1)$ \\
    Shaping velocità minima (gated$^{\dagger}$) & velocità sotto / sopra 8 km/h & $-0.04\,(8 - v)\cdot\text{align}$ / $+0.15\cdot\text{align}$ \\
    Rampa di non-progresso (gated) & $>$ 100 passi senza waypoint & $-0.004$ per passo, tetto $-1.0$ \\
    Sterzo fluido / a strattoni & $|\Delta \text{steer}| < 0.1$ a $> 5$ km/h \ / \ $> 0.5$ & $+0.1$ / $-0.3$ \\
    Anti-loop             & rilevatore di loop attivo & $-1.0$ \\
    \midrule
    \multicolumn{3}{@{}l}{\emph{Pedone}} \\
    Waypoint raggiunto / saltato & per waypoint avanzato / saltato & $+50$ / $-10$ \\
    Progresso denso       & variazione della distanza dal waypoint & $+2.0$ per metro guadagnato \\
    Collisione            & primo contatto (terminale) & $-50$ \\
    Marciapiede / carreggiata & per passo su marciapiede / corsia & $+0.2$ / $-0.3$ \\
    Passo di camminata    & velocità $\in [1.2, 2.6]$ m/s & $+0.3$ \\
    Inattività / corsa    & velocità $< 0.15$ m/s \ / \ $> 3.0$ m/s & $-0.15$ / $-0.7$ per m/s oltre \\
    Anti-loop             & rilevatore di loop attivo & $-1.0$ \\
    \bottomrule
  \end{tabular}
\end{table}

\chapter{Registro condensato degli esperimenti}
\label{app:registro}

La \cref{tab:candidati} elenca i candidati valutati durante la fase di
sviluppo guidata da gate (\cref{sec:gate}) con il verdetto e l'esito di
sintesi. Le run A/B sono da $3\times10^{5}$ passi, seed 999, bloccate su
easy salvo diversa indicazione; i delta sono rispetto alla baseline
appaiata della loro era e \emph{non} sono confrontabili tra versioni
dell'ambiente (\cref{sec:validita}). Il registro completo, con
meccanismi e identificatori di run per candidato, è mantenuto nel
repository.

\begin{table}[!htbp]
  \centering\small
  \caption{Registro condensato dei candidati.}
  \label{tab:candidati}
  \begin{tabular}{@{}L{4.5cm}L{2.6cm}L{2.6cm}L{4.3cm}@{}}
    \toprule
    Candidato & Tipo & Verdetto & Sintesi \\
    \midrule
    Correzione lunghezza percorsi & bug fix & adottato & $\approx +50$\,pp SR (cambio di distribuzione) \\
    Correzione seed percorsi & bug fix & adottato & ripristina l'appaiamento A/B \\
    Osservabilità 47D & osservazione & adottato & baseline sana, var.\ spiegata 0.983 \\
    Blocco di reward shaping anti-stallo & reward & promosso & nucleo della \cref{sec:reward} \\
    Rimozione guardia di progresso & reward & promosso & SR $+4.7$\,pp, s$+$t $-3.5$\,pp \\
    Tolleranza hazard $0.75 \to 0.85$ & reward & promosso & 5/5: SR $+2.3$\,pp, s$+$t $-3.0$\,pp \\
    Tetto attraversamenti $8 \to 3$ & percorsi & promosso & SR $+9.2$\,pp, coll.\ $-2.3$\,pp, off.\ $-2.2$\,pp \\
    Schedule di entropia a frazione & ottimizzatore & promosso & meccanismo stabile, nessuna instabilità \\
    Riparazione value loss & ottimizzatore & mantenuto (gate fallito) & var.\ spiegata $0 \to 0.87$; collisione $+3.3$\,pp \\
    Sconto $0.99 \to 0.997$ & ottimizzatore & mantenuto (falsificato) & timeout $\to$ collisione $\approx$ 1:1 \\
    Penalità di collisione $\times 10$ & reward & mantenuto (falsificato) & collisione piatta: limite di percezione, non di peso \\
    Terminazione precoce da stuck & dinamica env & rigettato & SR $-2.9$\,pp, s$+$t $+7.1$\,pp \\
    Retuning banda di passo pedoni & reward & rigettato & overshoot peggiore; declino in Q4 \\
    \texttt{route\_short} come successo & misura & revertito & violava il criterio di successo \\
    \bottomrule
  \end{tabular}
\end{table}

\chapter{Configurazione e iperparametri}
\label{app:config}

La \cref{tab:iperparametri} riporta la configurazione completa delle
quattro run di campagna. Tutti i valori sono identici tra le celle a
meno dei due flag sperimentali (\texttt{mode} e \texttt{use\_gnn}),
proprietà verificata con un diff diretto dei \texttt{run\_config.json}
scritti al lancio, che restano il riferimento autoritativo per ogni
singola run. Le famiglie di livelli per densità di traffico e miste sono
definite in configurazione ma inutilizzate in questa tesi
(\cref{sec:campagna}).

\begin{table}[!htbp]
  \centering\small
  \setlength{\tabcolsep}{4pt}
  \caption{Configurazione delle run di campagna.}
  \label{tab:iperparametri}
  \begin{tabular}{@{}llL{9.0cm}@{}}
    \toprule
    Gruppo & Parametro & Valore \\
    \midrule
    \multicolumn{3}{@{}l}{\emph{Ottimizzazione (PPO/MAPPO)}} \\
    & learning rate & $5\times10^{-4}$ \\
    & $\gamma$ & 0.997 \\
    & GAE $\lambda$ & 0.95 \\
    & clip PPO $\varepsilon$ & 0.2 \\
    & penalità KL & attiva; target 0.02, coefficiente iniziale 0.3 \\
    & gradient clipping & 0.5 \\
    & \texttt{vf\_clip\_param} & $10^{6}$ (clamp non vincolante) \\
    & \texttt{vf\_loss\_coeff} & 0.05 \\
    & coefficiente di entropia & schedule a frazione del budget: $0.03$ a $0.0$, $0.005$ a $0.83$ \\
    \addlinespace
    \multicolumn{3}{@{}l}{\emph{Rollout e batching}} \\
    & frammento di rollout & 200 passi \\
    & batch di addestramento & 8{.}000 passi \\
    & minibatch SGD & 256 \\
    & epoche SGD per batch & 15 \\
    & budget totale & $3\times10^{6}$ passi d'ambiente \\
    & frequenza di checkpoint & ogni $10^{5}$ passi \\
    \addlinespace
    \multicolumn{3}{@{}l}{\emph{Modello}} \\
    & actor (entrambe le classi) & MLP, 2 strati nascosti da 256 \\
    & critic MLP & $225 \to 256 \to 256 \to 1$ \\
    & critic GNN & embedding di nodo 64, 2 layer di message passing, aggregazione a media, ELU \\
    & PopArt, attenzione & disabilitati \\
    \addlinespace
    \multicolumn{3}{@{}l}{\emph{Curriculum}} \\
    & finestra di competenza & 50 episodi (deve essere piena) \\
    & metrica di sblocco & minimo dell'SR a finestra sulle due policy \\
    & soglia medium / hard & 0.70 / 0.60 \\
    & quota minima di budget & easy $\geq 0.20$ prima di medium; medium $\geq 0.28$ prima di hard \\
    & cap di pressione & medium al 30\% del budget, hard al 70\% \\
    & pesi di campionamento & $1.00 / 1.17 / 1.17$ ($\to 0.30/0.35/0.35$) \\
    & vincoli di budget & easy $\leq 0.30$, medium $\leq 0.35$, hard $\geq 0.35$ \\
    & probation & 2 blocchi dopo lo sblocco (1 dopo un cap); peso hard $0.85$ \\
    \addlinespace
    \multicolumn{3}{@{}l}{\emph{Batch}} \\
    & campionamento & mescolamento stratificato senza reinserimento, finestra 3 \\
    \addlinespace
    \multicolumn{3}{@{}l}{\emph{Livelli (difficoltà = lunghezza percorso)}} \\
    & easy / medium / hard & Town03, \SI{30}{m} / \SI{60}{m} / \SI{100}{m} \\
    & test (solo valutazione) & Town05, \SI{80}{m} \\
    & traffico NPC & 15 veicoli, 30 pedoni in ogni livello \\
    \bottomrule
  \end{tabular}
\end{table}

\chapter{Riproducibilità}
\label{app:riproducibilita}

\paragraph{Stack software.}
CARLA 0.9.16 (server e API Python); Python 3.11.9; PyTorch 2.7
(CUDA 12.6); Ray/RLlib 2.10.0. Addestramento e simulatore girano sulla
stessa macchina.

\paragraph{Hardware.}
Tutte le run sono state eseguite su una singola macchina di classe
laptop: Intel Core i7-10870H (8 core / 16 thread), \SI{16}{GB} di RAM,
NVIDIA GeForce RTX 3080 Laptop GPU (\SI{8}{GB}), Windows 11.
L'inviluppo di risorse è parte dei vincoli di design dello studio:
motiva il budget diagnostico da $3\times10^{5}$ passi del protocollo di
gate (\cref{sec:gate}) e la riduzione della campagna a un singolo seed
per cella (\cref{sec:campagna}).

\paragraph{Versione del codice e dati grezzi.}
Il trunk di campagna è il commit \texttt{7defb03}; i commit successivi
fino alla stesura di questa tesi toccano solo documentazione e tooling
di visualizzazione. Ogni directory di run conserva il proprio
\texttt{run\_config.json}, i checkpoint e \texttt{episodes.jsonl} --- i
record grezzi a livello di agente da cui ogni numero di questa tesi è
ricalcolato. Le quattro run di campagna sono
\texttt{20260622\_171626} (MLP-curriculum),
\texttt{20260623\_171855} (MLP-batch),
\texttt{20260630\_181143} (GNN-curriculum) e
\texttt{20260714\_155209} (GNN-batch).

\paragraph{Comandi.}
Addestramento (uno per cella; le quattro celle variano solo
\texttt{--mode} e \texttt{--use-gnn}):
\begin{quote}\small\ttfamily
python -m carla\_core.training.train\_carla\_mappo --mode curriculum\\
\hspace*{2em}--difficulty path --timesteps 3000000 --seed 999 [--use-gnn]
\end{quote}
L'addestramento scrive un \texttt{final\_eval\_job.json} nella directory
della run; la valutazione finale a 400 episodi si lancia con:
\begin{quote}\small\ttfamily
python -m carla\_core.training.evaluate\_carla\_mappo\\
\hspace*{2em}--job <run\_dir>/final\_eval\_job.json
\end{quote}

\paragraph{Seed.}
Campagna: un unico seed di esperimento (999), appaiato tra le quattro
celle; i seed dei percorsi per agente ne derivano deterministicamente
(\cref{sec:mondo}). Si noti che il seeding \emph{non} rende identiche le
run CARLA (\cref{sec:minacce}): riproducibilità qui significa riprodurre
la \emph{distribuzione} degli esiti, e ogni aggregato riportato a
partire dai log archiviati.
